\clearpage
\section*{NeurIPS Paper Checklist}

\begin{enumerate}

\item {\bf Claims}
\item[] Question: Do the main claims made in the abstract and introduction accurately reflect the paper's contributions and scope?
\item[] Answer: \answerYes{}
\item[] Justification: The abstract and \cref{sec:intro} state three contributions: (i) the \GeoGen\ methodology, (ii) the 801-prompt \GeoGenBench\ instance with a 17-predicate language, and (iii) headline empirical findings from a $600$-scenario $\times$ $3$-repeat draw across six frontier models and two strategies. All three are substantiated in the paper: (i) in \cref{sec:design}, (ii) in \cref{sec:design,app:property_reference}, and (iii) in \cref{sec:results,app:case_studies}. The Pareto-corner finding ($\gpt + \structured$ at $94.1$\%, \dollars{0.075}/scenario), the $7\times$ cost-equivalent-quality spread, the $\sim 18\%$ one-directional verifier lenience from the human-correlation study, and the Gemini-Flash deployment-portability finding are all reproducible from the released JSONL and tooling.

\item {\bf Limitations}
\item[] Question: Does the paper discuss the limitations of the work performed by the authors?
\item[] Answer: \answerYes{}
\item[] Justification: \cref{sec:discussion} (``Scope and limitations'' and ``Threats to validity'' paragraphs) enumerates the limitations of the benchmark: 2D Euclidean only, no proof construction, marginal partial coverage on two \structured cells, fixed tolerance $\tau$, $35\%$ of templated scenarios are loose, and the $\sim 18\%$ naturally-occurring verifier lenience. \cref{app:human_study:limitations} adds limitations specific to the human-correlation study (two raters, both paper authors; templated split only; three-day annotation window).

\item {\bf Theory assumptions and proofs}
\item[] Question: For each theoretical result, does the paper provide the full set of assumptions and a complete (and correct) proof?
\item[] Answer: \answerNA{}
\item[] Justification: The paper presents an empirical benchmarking study; it does not advance any theoretical results requiring formal proofs.

\item {\bf Experimental result reproducibility}
\item[] Question: Does the paper fully disclose all the information needed to reproduce the main experimental results of the paper to the extent that it affects the main claims and/or conclusions of the paper (regardless of whether the code and data are provided or not)?
\item[] Answer: \answerYes{}
\item[] Justification: \cref{app:reproduction} provides a complete reproduction recipe: the procedural template engine and exact seed (42), the per-(model, strategy, scenario, repeat) JSONL artefacts, the verification toolkit, the Docker renderer image, the API model identifiers and dates of access, and the scripts for re-aggregating the leaderboard. The hardening pass and its per-pilot-record promotion delta are documented in \cref{ssec:ver_defn}.

\item {\bf Open access to data and code}
\item[] Question: Does the paper provide open access to the data and code, with sufficient instructions to faithfully reproduce the main experimental results, as described in supplemental material?
\item[] Answer: \answerYes{}
\item[] Justification: The dataset (\texttt{benchmark/definitions/bench\_generated.jsonl}, \texttt{bench\_curriculum.jsonl}, plus a Croissant 1.0 metadata file), the procedural template engine, the verification toolkit, the rating viewer used in the human-correlation study, and the Docker renderer are all released under permissive licenses (MIT for code; details in \cref{app:reproduction} and the \texttt{benchmark/definitions/README.md} dataset card). Reviewers can re-roll the engine with a different seed for a contamination-free split.

\item {\bf Experimental setting/details}
\item[] Question: Does the paper specify all the training and test details (e.g., data splits, hyperparameters, how they were chosen, type of optimizer, etc.) necessary to understand the results?
\item[] Answer: \answerYes{}
\item[] Justification: \cref{sec:strategies} specifies model identifiers, sampling temperature ($1.0$, default), the two strategies' system prompts and tool-call interface, the $N{=}3$ self-correct retry budget for \structured, and the 600-scenario $\times$ 3-repeat headline-run shape. \cref{app:reproduction} adds the per-cell wall-clock and token totals and the subprocess-isolated-runner timeout configuration. We do not train any model in this work; we benchmark frozen production endpoints.

\item {\bf Experiment statistical significance}
\item[] Question: Does the paper report error bars suitably and correctly defined or other appropriate information about the statistical significance of the experiments?
\item[] Answer: \answerYes{}
\item[] Justification: The human-correlation primary endpoint reports $\kappa = 0.455$ with a $1000$-resample bootstrap $95\%$ CI $[0.304, 0.590]$ (\cref{ssec:human}, \cref{app:human_study:results}). Headline pass-rate spreads ($1$\pp on the three frontier+\structured\ cells, $20$\pp on the $\rawcode$ cross-vendor comparison) are at $N \geq 1{,}140$ records per cell, where Wilson $95\%$ CIs on a single cell are $\leq 1$\pp; cross-cell deltas of $\geq 5$\pp are well above sampling noise. Per-cell repeat stability is reported in \cref{app:case_studies:stability} ($\geq 96\%$ all-same-verdict on the three frontier+\structured\ cells), bounding within-cell variance.

\item {\bf Experiments compute resources}
\item[] Question: For each experiment, does the paper provide sufficient information on the computer resources (type of compute workers, memory, time of execution) needed to reproduce the experiments?
\item[] Answer: \answerYes{}
\item[] Justification: The headline run consumed approximately $\$975$ in API spend, with per-cell per-scenario cost reported in \cref{tab:headline} and per-cell wall-clock latency in the same table. Reproduction does not require GPUs or specialised hardware --- only API access to the named endpoints and a Docker host able to run the LuaLaTeX renderer container ($\sim 200$\,MB image, $<1$\,s/render). The eval harness uses asyncio with subprocess isolation and runs comfortably on a 4-core laptop. \cref{app:reproduction} documents the canonical compute environment.

\item {\bf Code of ethics}
\item[] Question: Does the research conducted in the paper conform, in every respect, with the NeurIPS Code of Ethics?
\item[] Answer: \answerYes{}
\item[] Justification: The work benchmarks publicly-available model APIs against a procedurally-generated dataset and a curriculum-derived dataset; it does not use personal data, does not deploy a system in a high-stakes setting, and does not introduce a generative model with elevated misuse risk. The human-correlation study uses two paper-author raters under documented conflict-of-interest mitigations (\cref{ssec:human}, \cref{app:human_study}).

\item {\bf Broader impacts}
\item[] Question: Does the paper discuss both potential positive societal impacts and negative societal impacts of the work performed?
\item[] Answer: \answerYes{}
\item[] Justification: \emph{Positive impacts:} an automatic, geometry-aware verifier for diagram generation enables faster iteration on LLM-driven educational content (e.g., K--12 worksheet authoring, technical illustrations) and provides a cheap reward signal for fine-tuning that does not require human labels or vision-language judges. The procedural template engine ships with the benchmark so reviewers and downstream researchers can re-roll contamination-free splits, lowering the cost of evaluation hygiene. \emph{Negative impacts:} an automatically-graded benchmark could displace human assessment in educational technology where teaching utility (rather than dimensional correctness) is the relevant signal --- we explicitly note in \cref{sec:discussion} that \GeoGenBench\ does not measure teaching utility and should not be used as a proxy for it. The human-correlation study's $18\%$ one-directional lenience finding is also published transparently to caution downstream users about over-relying on the auto-verdict for student-facing decisions.

\item {\bf Safeguards}
\item[] Question: Does the paper describe safeguards that have been put in place for responsible release of data or models that have a high risk for misuse?
\item[] Answer: \answerNA{}
\item[] Justification: The paper releases a benchmark dataset, a verification toolkit, and a Docker renderer; it does not release a generative model or scraped dataset. The released artefacts are low-risk: the dataset consists of procedurally-generated geometry prompts and machine-decidable property lists; misuse risk is comparable to releasing a unit-test suite. The curriculum-extracted scenarios cite their topical scaffold (Carnegie Learning's \emph{Bluebonnet Learning} curriculum, an open public-domain curriculum) and reproduce no copyrighted prose.

\item {\bf Licenses for existing assets}
\item[] Question: Are the creators or original owners of assets used in the paper properly credited and are the license and terms of use explicitly mentioned and properly respected?
\item[] Answer: \answerYes{}
\item[] Justification: The curriculum scaffold is the Texas Education Agency / Carnegie Learning's \emph{Bluebonnet Learning} K--12 curriculum, attributed in \cref{sec:design,app:reproduction} and \texttt{docs/datasheet.md}; we use Bluebonnet only as a topical scaffold (chapter / lesson structure) and reproduce no Bluebonnet prose, figures, or problems. The TikZ ecosystem (\texttt{tkz-euclide}, \texttt{tkz-elements}, \texttt{tikz}) is licensed under the LaTeX Project Public License; SymPy under BSD-3-Clause; \texttt{pydantic-ai} under MIT; the model APIs (Anthropic, OpenAI, Google) under their respective developer terms. All upstream licenses are listed in \cref{app:reproduction} and \texttt{benchmark/definitions/README.md}.

\item {\bf New assets}
\item[] Question: Are new assets introduced in the paper well documented and is the documentation provided alongside the assets?
\item[] Answer: \answerYes{}
\item[] Justification: The released benchmark ships with: a Croissant 1.0 metadata file (\texttt{benchmark/definitions/croissant.json}, validated against \texttt{mlcroissant} 1.1.0), a Hugging Face-style dataset card (\texttt{benchmark/definitions/README.md}), a Gebru-format datasheet (\texttt{docs/datasheet.md}), a frozen pre-registration for the human-correlation study (\texttt{docs/human\_study\_protocol.md}), and a public leaderboard URL committed to the repository.

\item {\bf Crowdsourcing and research with human subjects}
\item[] Question: For crowdsourcing experiments and research with human subjects, does the paper include the full text of instructions given to participants and screenshots, if applicable, as well as details about compensation?
\item[] Answer: \answerYes{}
\item[] Justification: \cref{app:human_study:rubric} reproduces the full rating rubric (Q1, Q2, blinding); \cref{app:human_study:tooling} describes the rating viewer (single-page HTML, screenshot deferred to the released repository at \texttt{evals/human\_study/viewer.html} and the runbook at \texttt{evals/human\_study/README.md}); the protocol is committed at \texttt{docs/human\_study\_protocol.md} prior to data collection. Both raters are paper authors; no monetary compensation was given; the conflict-of-interest mitigation (Q1 blinding + per-rater independent $\kappa$) is reported in \cref{ssec:human,app:human_study:limitations}.

\item {\bf Institutional review board (IRB) approvals or equivalent}
\item[] Question: Does the paper describe potential risks incurred by study participants, whether such risks were disclosed to the subjects, and whether IRB approvals (or an equivalent approval/review based on the requirements of your country or institution) were obtained?
\item[] Answer: \answerNA{}
\item[] Justification: The human-correlation study involves only the two paper authors as raters, evaluates publicly-available LLM outputs against publicly-released geometry prompts, collects no personal information beyond a self-chosen rater identifier, and presents no foreseeable risk to the raters beyond the time investment of $7$--$13$ hours per rater of focused rating work. Under our institution's review process, low-risk self-rating studies of this kind do not require IRB approval. The annotation effort and rater identities are disclosed in \cref{ssec:human,app:human_study:raters}.

\item {\bf Declaration of LLM usage}
\item[] Question: Does the paper describe the usage of LLMs if it is an important, original, or non-standard component of the core methods in this research?
\item[] Answer: \answerYes{}
\item[] Justification: LLMs play two distinct roles in this work, both disclosed: (i) the systems under evaluation --- six production frontier models from three vendors (Anthropic Claude Opus 4.7 / Sonnet 4.6 / Haiku 4.5; OpenAI GPT-5.5; Google Gemini 2.5 Pro / Flash) are benchmarked end-to-end (\cref{sec:strategies}); and (ii) the curriculum-split scenario authoring --- Anthropic Claude Sonnet 4.6 was used to convert lesson-topic scaffolds from the \emph{Bluebonnet Learning} curriculum into student-style diagram-drawing prompts paired with author-written property lists (\cref{sec:design,app:reproduction}), with the templated split (procedurally-generated, no LLM in the authoring loop) reported separately as a calibration anchor.

\end{enumerate}
